@@M0@@ which sends @@M1@@ to @@M2@@, and @@M3@@ to @@M4@@; by the preceding remark, the element obtained does not depend on the integer @@M5@@ such that @@M6@@ for @@M7@@. Using the same homomorphism, we see that relation (5.3.1) is satisfied by @@M8@@ and @@M9@@.

\sgaStableTarget{sga6.e05.para.5-4}5.4. Consider the @@M10@@-ring generated over @@M11@@ by the indeterminates @@M12@@, subject to the relations @@M13@@ for @@M14@@. From (5.3.1) we obtain:
\begin{align*}
  \lambda^p(N,x+y)\mathbin{.}\lambda_{-1}(N)
  &=\lambda^p\bigl((x+y)\mathbin{.}\lambda_{-1}(N)\bigr)\\
  &=\sum_{i+j=p}\lambda^i(x\mathbin{.}\lambda_{-1}(N))\mathbin{.}\lambda^j(y\mathbin{.}\lambda_{-1}(N))\\
  &=\sum_{\substack{i+j=p\\ i,j\ne0}}\lambda^i(N,x)\mathbin{.}\lambda^j(N,y)\mathbin{.}(\lambda_{-1}(N))^2
    +(\lambda^p(N,x)+\lambda^p(N,y))\lambda_{-1}N
\end{align*}

Since the ring under consideration is an integral domain, we deduce the following relation, valid by substitution in every @@M15@@-ring, and for every element @@M16@@ such that @@M17@@ for @@M18@@:
\sgaStableTarget{sga6.e05.eq.5-4-1}@@M19@@

Now consider the @@M20@@-ring generated over @@M21@@ by indeterminates @@M22@@ subject to the relations @@M23@@ for @@M24@@, and @@M25@@ for @@M26@@. We then have @@M27@@ for @@M28@@; we can therefore define the @@M29@@. We obtain:
@@M30@@
By the additivity of @@M31@@, we have: @@M32@@. We deduce:
\begin{align*}
  \lambda^p(N+N',x)\mathbin{.}\lambda_{-1}(N)\mathbin{.}\lambda_{-1}(N')
  &=\lambda^p(x\mathbin{.}\lambda_{-1}(N)\mathbin{.}\lambda_{-1}(N'))\\
  &=\lambda^p(N',x\mathbin{.}\lambda_{-1}(N))\mathbin{.}\lambda_{-1}(N').
\end{align*}

